% =====================================================================
% sections/empirical.tex — shared (long + short)
% =====================================================================

\section{Empirical strategy}
\label{sec:empirical}

\subsection{The identification problem}

The Village Midwife Program was not rolled out at random.
Provincial health departments deliberately prioritized poor and remote
communities with the weakest existing health infrastructure when
allocating midwife graduates \autocite{frankenberg2005,gani1996,
makowiecka2008}. As a result, the year a community first received a
village midwife, $\text{start}_c$, is mechanically correlated with
village-level characteristics---baseline poverty, distance to the
nearest clinic, ethnic composition---that themselves affect child
outcomes through channels other than the program. A naïve comparison
of cohorts treated early to those treated late conflates the program's
effect with the consequences of these confounders. Inverse-probability
weighting on observable village characteristics
\autocite{cas2012}, the approach used in the predecessor DEPII analysis,
is one defense, but it cannot correct for unobserved village
attributes that drove placement decisions.

A second concern, distinct from placement bias, is that the canonical
two-way fixed-effects estimator is known to be biased in the presence
of staggered treatment timing with heterogeneous effects, because it
implicitly uses already-treated units as controls for later-treated
ones \autocite{goodmanbacon2021}. \textcite{ahsan2021} estimate the
program's effects using a single OLS regression with community fixed
effects absorbed; their framework does not separate cohort-level
heterogeneity from staggered-treatment bias.

The empirical design responds to both concerns with a layered
strategy. The headline TWFE estimate is reported alongside the
\textcite{callawaysantanna2021} estimator, which uses never-treated
communities as a clean control group, and the
\textcite{dechaisemartin2021} estimator. Pre-trends are tested by an
event-study following the Sun--Abraham aggregation. The placement-bias
problem is addressed via a within-mother specification that sweeps out
all time-invariant village characteristics that drove rollout order.
\textcite{anderson2008}-style composite indices reduce multiple-testing
inflation across closely related outcomes within a domain. Finally,
\textcite{frankenberg2005}'s observation that the rollout was
demand-driven by community characteristics motivates the
\textcite{ahsan2021} sensitivity tests that I revisit in
Section~\ref{sec:robustness}.

\subsection{Baseline two-way fixed-effects specification}

Let $i$ index individuals, $c$ their community of birth, and $b$ their
year of birth. The baseline specification is
\begin{equation}
Y_{ic} \;=\; \beta\, E_{ic} \;+\; X_{ic}'\theta \;+\; \eta_c \;+\;
\gamma_b \;+\; \varepsilon_{ic},
\label{eq:twfe}
\end{equation}
where $Y_{ic}$ is one of the four standardized outcome indices, $E_{ic}$
is an indicator equal to one if individual $i$ was in utero or aged
0--3 at the time the village midwife arrived in $c$,
$X_{ic}$ is a vector of pre-determined controls (sex, mother's
years of schooling, mother's age at birth), $\eta_c$ and $\gamma_b$ are
community-of-birth and birth-year fixed effects, and standard errors
are clustered at the community-of-birth level (\NVillages~clusters).
The age-three threshold is motivated by the developmental window
literature: the prenatal period and the first three years of life are
when nutrition, infection control, and preventive care have the
largest documented long-run consequences \autocite{barker1995}.

Two extensions of equation~\eqref{eq:twfe} are reported in the main
results table. The first adds the pre-determined controls listed above.
The second additionally absorbs province-by-birth-year fixed effects,
which allow each province to have its own arbitrary cohort trend and
guards against confounding from province-specific shocks.

\subsection{Heterogeneity-robust estimators}

The \textcite{callawaysantanna2021} estimator computes
group-time average treatment effects $ATT(g,t)$ for each
cohort group $g$ and each calendar period $t$, using never-treated
communities as the control group, and then aggregates them. I report
the dynamic aggregation
\[
ATT^{\text{dyn}}(e) \;=\; \frac{1}{|G_e|}\sum_{g \in G_e}
ATT\bigl(g,\, g+e\bigr),
\]
averaged over event times $e = 0, 1, 2, 3$, to match the same
``in-utero or age 0--3'' window as $E_{ic}$. The
\textcite{dechaisemartin2021} estimator follows a similar logic
but with a different aggregation scheme; both serve as a check that
the headline TWFE estimate is not driven by forbidden comparisons of
the kind highlighted by \textcite{goodmanbacon2021}.

\subsection{Event study}

The Sun--Abraham specification estimates a cohort-aggregated
event-study version of equation~\eqref{eq:twfe} via
\texttt{fixest::sunab}, with $-1$ as the omitted reference event time:
\begin{equation}
Y_{ic} \;=\; \sum_{e \neq -1} \delta_e \,\mathbf{1}\{b_i - g_c = e\}
\;+\; \eta_c \;+\; \gamma_b \;+\; u_{ic}.
\label{eq:sunab}
\end{equation}
A joint Wald test of $\delta_{-4} = \delta_{-3} = \delta_{-2} = 0$
provides a parallel-trends diagnostic. The displayed event-time window
is $[-4, +8]$; cohorts more than four years before rollout are folded
into the omitted reference because their identification rests on a
small number of villages whose rollout was very late.

\subsection{Sibling fixed effects: the headline robustness}

If placement bias operates through time-invariant village
characteristics---poverty, remoteness, baseline health
infrastructure---then comparing siblings born to the same mother in
the same village but with different exposure to the village midwife
sweeps out every such confounder, observable or unobservable. The
within-mother specification is
\begin{equation}
Y_{ic} \;=\; \beta\, E_{ic} \;+\; X_{ic}'\theta \;+\; \alpha_{m(i)}
\;+\; \gamma_b \;+\; \varepsilon_{ic},
\label{eq:sibfe}
\end{equation}
where $\alpha_{m(i)}$ is a fixed effect for the mother of $i$, $X_{ic}$
adds birth order and mother's age at birth (siblings differ on these
even within mother), and standard errors remain clustered at the
community-of-birth level. Identification comes from differential
exposure across siblings born to the same mother before versus after
the midwife's arrival. The DEPII version of this study (Section~6,
``Extensions'') flagged sibling fixed effects as a natural next step
that the analysis frame supports; the present pipeline observes
\num{2796} mothers with two or more children in the analysis sample,
of whom \num{1153} have three or more.

What sibling fixed effects do \emph{not} solve is also worth stating
explicitly. Time-varying village shocks correlated with rollout
year---for instance, the simultaneous arrival of a midwife and a
paved road---will load on $\beta$ as if they were the program. The
empirical strategy addresses this by including controls for the
arrival of other health facilities (Section~\ref{sec:robustness},
column~R2). Compositional shifts in fertility induced by the midwife's
arrival---for example, if mothers time births around the program---are
a separate threat that I check in Section~\ref{sec:robustness}.

\subsection{Inverse-probability weighting (demoted)}

The DEPII version of this analysis used inverse-probability weighting
on baseline community characteristics as the principal correction for
non-random rollout. In light of the layered strategy above I demote
IPW to a single robustness column (Section~\ref{sec:robustness},
column~R1). Weights are constructed from a logistic propensity score
of an early-treatment indicator on province dummies and post-period
other-facility presence, then trimmed at the 1st and 99th percentiles
and rescaled so $\bar{w} = 1$. The mechanical weakness of IPW in this
setting---it cannot correct for selection on unobservables---is one of
the motivations for adopting sibling fixed effects as the primary
robustness check.

\subsection{Identifying assumptions}

The TWFE estimator in equation~\eqref{eq:twfe} requires a
parallel-trends assumption: in the absence of the program, treated and
control cohorts within a community would have evolved similarly.
The Callaway--Sant'Anna estimator weakens this to parallel trends
between each cohort group and the never-treated, with no
extrapolation across cohort groups. Sibling fixed effects in
equation~\eqref{eq:sibfe} additionally require that the variation in
sibling exposure within mother is exogenous to the outcome conditional
on $\alpha_m$, $\gamma_b$, birth order, and maternal age---roughly,
that mothers do not time births in anticipation of the midwife's
arrival in a way that correlates with child outcomes through channels
other than midwife exposure. The fertility-selection check in
Section~\ref{sec:robustness} examines this. SUTVA requires that
treatment of one community does not affect outcomes in another, which
\textcite{ahsan2021} discuss in the context of midwife mobility across
village boundaries.

The balance test reported in Table~\ref{tab:balance} regresses
pre-program village covariates on early-treatment status and confirms
that placement was not random---a feature, not a bug: the test exists
to document the very confounding the rest of the empirical strategy
seeks to address.

\begin{table}[!htbp]
\centering
\caption{Balance: pre-program village covariates by early- vs.\
late-treatment status.}
\label{tab:balance}
% Auto-generated by code/R/07_balance_descriptives.R
\begin{tabular}{lrrrr}
\toprule
Variable & Early-rollout mean & Late-rollout mean & Difference & $p$-value \\
\midrule
Mother's education (years) & 5.154 & 6.383 & -1.228$^{***}$ & 0.000 \\
Mother's age at birth (years) & 26.441 & 27.489 & -1.049$^{***}$ & 0.000 \\
Mother's height (cm) & 150.056 & 150.298 & -0.242 & 0.248 \\
Female & 0.505 & 0.504 & 0.000 & 0.988 \\
\midrule
$N$ (children) & \multicolumn{2}{c}{2,017 / 1,955} & & \\
\bottomrule
\end{tabular}

\begin{minipage}{0.85\textwidth}
\footnotesize
\emph{Notes.} Sample is at the community-of-birth level
($\NVillages$ communities). Early-treated communities are those whose
SAR rollout year is at or below the median across treated communities.
$p$-values are from two-sample $t$-tests.
\end{minipage}
\end{table}
